\documentclass[11pt,letterpaper]{article}

\usepackage[margin=1in]{geometry}
\usepackage{amsmath,amssymb}
\usepackage{graphicx}
\usepackage{booktabs}
\usepackage{longtable}
\usepackage{array}
\usepackage{xcolor}
\usepackage{hyperref}
\usepackage{float}
\usepackage{enumitem}
\usepackage{caption}

\hypersetup{
    colorlinks=true,
    linkcolor=blue!55!black,
    urlcolor=blue!55!black,
    citecolor=blue!55!black
}

\graphicspath{{../../outputs/}}

\title{\textbf{Energy-Gated Frequency Neuron}\\
\large Un bloque neuronal fisico-inspirado para clasificacion de audio}
\author{Proyecto de portafolio e investigacion aplicada}
\date{Julio 2026}

\begin{document}
\maketitle

\begin{abstract}
Este documento describe el desarrollo conceptual, matematico y experimental de
la \textit{Energy-Gated Frequency Neuron} (EGFN), un bloque neuronal para audio
que combina bancos de filtros, activaciones no lineales, energia por banda de
frecuencia y compuertas aprendibles. El objetivo del proyecto no es demostrar
un nuevo estado del arte, sino construir una pieza de investigacion aplicada
que conecte procesamiento de senales, redes neuronales e interpretabilidad. Se
presenta la motivacion inicial, la formulacion del bloque, la arquitectura
implementada en PyTorch, las variantes experimentales evaluadas y una
comparacion con enfoques alternativos como MFCC+MLP, Mel-spectrogram+CNN,
Conv1D, SincNet y LEAF.
\end{abstract}

\tableofcontents
\newpage

\section{Origen del proyecto}

El proyecto surge de una pregunta conceptual: \textit{los filtros que aparecen
en circuitos y procesamiento de senales, como filtros pasa baja, pasa alta y
pasa banda, pueden utilizarse como funciones de activacion en redes
neuronales?}

La respuesta corta es que no son equivalentes. Una funcion de activacion
clasica opera punto a punto:

\begin{equation}
    y_i = f(x_i).
\end{equation}

En cambio, un filtro digital opera sobre una vecindad temporal:

\begin{equation}
    y[n] = \sum_k h[k]x[n-k].
\end{equation}

Por lo tanto, un filtro introduce estructura temporal y memoria local, pero
sigue siendo una transformacion lineal si no se combina con una no linealidad.
Esta diferencia es importante: reemplazar ReLU por un filtro no produce
automaticamente una red profunda no lineal. La idea fuerte del proyecto es otra:
usar filtros como una primera etapa fisico-inspirada y luego combinarlos con
activaciones, energia y compuertas aprendibles.

La hipotesis de trabajo es:

\begin{quote}
Un modelo de audio puede ganar interpretabilidad y potencialmente robustez si
procesa la senal como una composicion de bandas de frecuencia, energia por
banda y decisiones aprendibles de paso o atenuacion.
\end{quote}

\section{Principios de senales aplicados}

\subsection{Audio como senal temporal}

Un audio mono discreto puede representarse como una secuencia:

\begin{equation}
    x[n], \quad n = 0,1,\ldots,T-1.
\end{equation}

En deep learning, esta senal normalmente se organiza como un tensor:

\begin{equation}
    x \in \mathbb{R}^{B \times 1 \times T},
\end{equation}

donde $B$ es el tamano de batch, $1$ es el canal de audio mono y $T$ es el
numero de muestras temporales.

\subsection{Filtros y convolucion}

Un filtro digital se define por un kernel $h[k]$. La salida filtrada es:

\begin{equation}
    u[n] = (h * x)[n] = \sum_k h[k]x[n-k].
\end{equation}

Esto conecta directamente con las redes convolucionales. Una capa
\texttt{Conv1D} tambien aprende kernels que recorren la senal. Por tanto, una
CNN 1D ya contiene una idea de filtrado, aunque sus filtros no necesariamente
tengan una interpretacion fisica directa.

\subsection{No linealidad}

Si una red contiene solo transformaciones lineales, la composicion completa
puede reducirse a una unica transformacion lineal:

\begin{equation}
    y = W_3(W_2(W_1x)) = W_{\text{total}}x.
\end{equation}

Por eso se necesita una activacion:

\begin{equation}
    y = W_3 f(W_2 f(W_1x)).
\end{equation}

En la EGFN, el filtro prepara la senal y la activacion introduce capacidad no
lineal.

\subsection{Energia por banda}

En procesamiento de senales, la energia mide la intensidad o actividad de una
senal. Para la respuesta filtrada $u_k[n]$ de una banda $k$, se usa:

\begin{equation}
    E_k = \frac{1}{T}\sum_n u_k[n]^2.
\end{equation}

Esta magnitud es interpretable: si $E_k$ es alta, la banda $k$ contiene
actividad importante para esa senal.

\section{Energy-Gated Frequency Neuron}

\subsection{Formulacion principal}

Para cada banda de frecuencia $k$, la EGFN calcula:

\begin{align}
    u_k[n] &= (h_k * x)[n], \\
    z_k[n] &= u_k[n] + b_k, \\
    a_k[n] &= \phi(z_k[n]), \\
    E_k &= \frac{1}{T}\sum_n u_k[n]^2, \\
    g_k &= \sigma(\alpha_k \log(1 + E_k) + \beta_k), \\
    y_k[n] &= g_k a_k[n].
\end{align}

Donde:

\begin{itemize}[leftmargin=*]
    \item $h_k$ es el filtro de la banda $k$.
    \item $\phi$ es una activacion, por ejemplo GELU o ReLU.
    \item $E_k$ es la energia de la banda.
    \item $g_k$ es una compuerta aprendible entre 0 y 1.
    \item $\alpha_k$, $\beta_k$ y $b_k$ son parametros entrenables.
\end{itemize}

\subsection{Interpretacion}

El bloque se puede leer como:

\begin{quote}
Filtra la senal por banda, activa la respuesta, mide cuanta energia contiene
esa banda y aprende cuanto dejar pasar hacia el clasificador.
\end{quote}

Esto permite visualizar no solo que clase predice el modelo, sino tambien que
bandas de frecuencia usa para llegar a esa decision.

\section{Arquitectura neuronal completa}

La red implementada tiene tres componentes principales:

\begin{enumerate}[leftmargin=*]
    \item \textbf{Frontend EGFN}: banco de filtros, activacion, energia y gate.
    \item \textbf{FrequencyPooling}: resumen temporal interpretable.
    \item \textbf{Classifier Head}: MLP pequeno para clasificacion.
\end{enumerate}

\subsection{Flujo de tensores}

\begin{verbatim}
Raw waveform [B, 1, T]
  -> EnergyGatedFrequencyNeuron
       -> filter bank
       -> activation
       -> energy estimation
       -> learned gate
  -> FrequencyPooling
       -> mean, max, std, log-energy, gates
  -> MLP classifier
  -> logits [B, num_classes]
\end{verbatim}

\subsection{Pooling temporal}

Despues de la EGFN, se resumen las respuestas por banda:

\begin{align}
    \mu_k &= \text{mean}_n(y_k[n]), \\
    m_k &= \max_n(y_k[n]), \\
    s_k &= \text{std}_n(y_k[n]).
\end{align}

Luego se concatenan:

\begin{equation}
    p = [\mu, m, s, \log(1+E), g].
\end{equation}

Esta representacion es compacta y mantiene informacion interpretable por banda.

\section{Variantes implementadas}

\subsection{V0: filtros fijos}

La primera version usa filtros predefinidos tipo pasa baja, pasa banda y pasa
alta. Esta version es la mas interpretable porque cada filtro conserva su
significado original.

\subsection{V1: filtros aprendibles}

En la segunda version, los filtros se inicializan como filtros interpretables,
pero sus kernels se vuelven parametros entrenables. Esto aumenta flexibilidad,
pero puede reducir interpretabilidad si los kernels dejan de parecer filtros
limpios.

\subsection{Banco de filtros fino}

El banco original tenia 5 bandas. Para el caso dificil, se agrego un banco fino
de 8 bandas:

\begin{center}
\begin{tabular}{ll}
\toprule
Indice & Banda \\
\midrule
1 & lowpass \\
2 & 250--500 Hz \\
3 & 500--750 Hz \\
4 & 750--1000 Hz \\
5 & 1000--1500 Hz \\
6 & 1500--2500 Hz \\
7 & 2500--4000 Hz \\
8 & highpass \\
\bottomrule
\end{tabular}
\end{center}

Este cambio fue importante porque las clases hard viven principalmente entre
aproximadamente 420 Hz y 1260 Hz.

\subsection{Gate contextual}

El gate independiente calcula:

\begin{equation}
    g_k = \sigma(\alpha_k \log(1 + E_k) + \beta_k).
\end{equation}

El gate contextual usa todo el vector de energias:

\begin{equation}
    g = \sigma(\text{MLP}([\log(1+E_1), \ldots, \log(1+E_K)])).
\end{equation}

La motivacion era permitir que cada banda se interprete en contexto. Sin
embargo, los resultados iniciales no mostraron mejora clara. Esta observacion
es valiosa: mas capacidad no siempre mejora el modelo.

\subsection{V2: sparse selective gate}

La V2 introduce una penalizacion:

\begin{equation}
    \mathcal{L} =
    \mathcal{L}_{CE} + \lambda_{gate}\frac{1}{BK}\sum_{b,k} g_{b,k}.
\end{equation}

La intencion es que el modelo no use todas las bandas si no es necesario. Los
resultados iniciales muestran selectividad suave: no se apagan todas las bandas
de forma binaria, pero con umbral 0.7 se observa que el modelo concentra gates
fuertes en cerca de dos bandas promedio.

\section{Datos sinteticos}

Antes de usar audio real, se construyo un dataset sintetico controlado. Esto es
importante porque permite saber si el bloque responde a frecuencia y energia
antes de introducir complejidad externa.

\subsection{Dificultad easy}

Las clases estan muy separadas:

\begin{itemize}[leftmargin=*]
    \item Clase 0: bajas frecuencias.
    \item Clase 1: frecuencias medias.
    \item Clase 2: frecuencias altas.
    \item Clase 3: mezcla baja + alta.
\end{itemize}

Este experimento valido que el pipeline completo funciona, pero resulto muy
facil.

\subsection{Dificultad hard}

El modo hard agrega:

\begin{itemize}[leftmargin=*]
    \item frecuencias mas cercanas,
    \item jitter de frecuencia,
    \item tonos distractores,
    \item ruido controlado por SNR,
    \item desplazamientos temporales,
    \item dropout parcial de componentes.
\end{itemize}

Este modo es mas util para investigar si el modelo realmente aprende
representaciones robustas.

\section{Resultados experimentales iniciales}

Los resultados aun son exploratorios y se basan en datos sinteticos. No deben
interpretarse como una comparacion definitiva contra modelos de audio reales.

\begin{longtable}{p{0.38\linewidth}rrrrr}
\caption{Resumen de experimentos locales.}\\
\toprule
Configuracion & Best epoch & Best val acc & Final val acc & Gate mean & Active bands \\
\midrule
\endfirsthead
\toprule
Configuracion & Best epoch & Best val acc & Final val acc & Gate mean & Active bands \\
\midrule
\endhead
Hard + fine + fixed filters & -- & -- & -- & -- & -- \\
Hard + fine + learnable filters & -- & -- & -- & -- & -- \\
Hard + learnable + contextual gate & 40--90 & 0.703 & 0.688 & -- & -- \\
Hard + sparse gate, $\lambda=0.005$ & -- & 0.711 & -- & 0.667 & 8.00 \\
Hard + sparse gate, $\lambda=0.05$ & 92 & 0.727 & 0.688 & 0.623 & 7.99 \\
Sparse gate, threshold 0.7 & 92 & 0.719 & 0.695 & 0.651 & 2.07 \\
\bottomrule
\end{longtable}

\subsection{Lectura de resultados}

\begin{itemize}[leftmargin=*]
    \item El modo easy alcanzo accuracy perfecta rapidamente, validando el
    pipeline pero mostrando que el problema era demasiado simple.
    \item El modo hard con banco grueso se quedo cerca de 0.53, lo que sugirio
    falta de resolucion en frecuencia.
    \item El banco fine y filtros aprendibles mejoraron el comportamiento en
    hard.
    \item El gate contextual no mejoro automaticamente. Es una ablation util.
    \item La sparse gate mejoro ligeramente accuracy y mostro selectividad
    suave.
\end{itemize}

\section{Figuras experimentales}

\begin{figure}[H]
\centering
\includegraphics[width=\linewidth]{hard_sparse_gate_l005/figures/synthetic_training_curves.png}
\caption{Curvas de entrenamiento para sparse gate con $\lambda=0.05$.}
\end{figure}

\begin{figure}[H]
\centering
\includegraphics[width=0.95\linewidth]{hard_sparse_gate_l005/figures/synthetic_gates_by_class.png}
\caption{Gates promedio por clase y banda.}
\end{figure}

\begin{figure}[H]
\centering
\includegraphics[width=0.95\linewidth]{hard_sparse_gate_l005/figures/synthetic_energy_by_class.png}
\caption{Energia logaritmica promedio por clase y banda.}
\end{figure}

\begin{figure}[H]
\centering
\includegraphics[width=\linewidth]{hard_sparse_gate/figures/snr_sweep_hard.png}
\caption{Barrido de robustez a ruido para el caso hard sparse gate.}
\end{figure}

\section{Comparacion con arquitecturas alternativas}

\subsection{MFCC + MLP}

MFCC es un pipeline clasico de audio. Extrae caracteristicas compactas basadas
en una escala perceptual y luego un MLP puede clasificarlas. Es eficiente y
fuerte como baseline, pero la extraccion de caracteristicas esta predefinida.
La EGFN, en cambio, aprende gates y opcionalmente filtros, manteniendo una
lectura por banda.

\subsection{Mel-spectrogram + CNN}

Este enfoque transforma audio en una representacion tiempo-frecuencia y aplica
convoluciones 2D. Es un baseline fuerte para audio. Su principal ventaja es
que produce una imagen espectral rica. Su desventaja para este proyecto es que
la interpretacion del frontend queda parcialmente fijada por el diseno mel.

\subsection{Conv1D sobre waveform}

Una CNN 1D aprende directamente desde la forma de onda. Es flexible, pero sus
primeros kernels pueden ser dificiles de interpretar. La EGFN busca un punto
medio: usar filtros y gates con significado fisico.

\subsection{SincNet}

SincNet aprende filtros pasa banda parametrizados por frecuencias de corte baja
y alta. Esto es muy relevante para la futura V3 del proyecto, porque permite
aprender filtros sin perder interpretabilidad fisica. A diferencia de kernels
libres, los parametros siguen teniendo significado.

\subsection{LEAF y EfficientLEAF}

LEAF propone un frontend aprendible completo para audio: filtros, pooling,
compresion y normalizacion. EfficientLEAF muestra que estas ideas pueden ser
mas eficientes, pero tambien advierte que los frontends aprendibles no siempre
superan a mel-filterbanks fijos. Esto refuerza la necesidad de baselines y
ablations honestas.

\section{Como se podria mejorar}

\subsection{Gate sharpness}

La sparse gate actual produce selectividad suave. Una mejora directa seria:

\begin{equation}
    g = \sigma(\tau \cdot r),
\end{equation}

donde $r$ es el logit del gate y $\tau > 1$ aumenta la nitidez. Esto podria
empujar gates hacia valores mas cercanos a 0 o 1.

\subsection{Filtros parametrizados tipo SincNet}

En vez de aprender todos los valores del kernel, se pueden aprender parametros:

\begin{equation}
    h_k = \text{BandPass}(f_{low,k}, f_{high,k}).
\end{equation}

Esto permitiria decir, por ejemplo, que una banda aprendio a moverse hacia
620--940 Hz. Para un proyecto de portafolio, esta mejora seria muy fuerte
porque une capacidad aprendible e interpretabilidad.

\subsection{Temporal Conv Head}

El pooling actual resume toda la senal con estadisticas globales. Para audio
real puede ser importante cuando ocurre la energia, no solo cuanta energia hay.
Una cabeza temporal podria usar:

\begin{verbatim}
EGFN output [B, K, T]
  -> temporal Conv1D
  -> pooling
  -> classifier
\end{verbatim}

Esto mantiene la EGFN como frontend interpretable y agrega modelado temporal.

\subsection{Datos reales}

Despues de validar con datos sinteticos, el siguiente paso es usar datasets
reales:

\begin{itemize}[leftmargin=*]
    \item Free Spoken Digit Dataset para una primera prueba liviana.
    \item Speech Commands para keyword spotting.
    \item ESC-50 o UrbanSound8K para eventos ambientales.
\end{itemize}

\section{Conclusiones}

La Energy-Gated Frequency Neuron es una arquitectura experimental razonable
para un portafolio de investigacion. Su valor no esta solo en buscar accuracy,
sino en mostrar una cadena completa de razonamiento:

\begin{itemize}[leftmargin=*]
    \item partir de filtros y activaciones,
    \item distinguir operaciones lineales y no lineales,
    \item disenar un bloque fisico-inspirado,
    \item implementar una capa personalizada en PyTorch,
    \item evaluar variantes y ablations,
    \item interpretar gates y energias por banda,
    \item comparar contra modelos existentes.
\end{itemize}

Los resultados actuales sugieren que filtros aprendibles y bancos de frecuencia
mas finos ayudan en problemas ambiguos. El gate contextual no mejoro
automaticamente, mientras que la sparse gate mostro selectividad suave y una
leve mejora de accuracy. La siguiente etapa mas prometedora es implementar
filtros parametrizados tipo SincNet o una cabeza temporal ligera.

\section*{Referencias}

\begin{itemize}[leftmargin=*]
    \item Ravanelli, M. and Bengio, Y. \textit{Speech and Speaker Recognition
    from Raw Waveform with SincNet}. arXiv:1812.05920.
    \url{https://arxiv.org/abs/1812.05920}
    \item Zeghidour, N. et al. \textit{LEAF: A Learnable Frontend for Audio
    Classification}. arXiv:2101.08596.
    \url{https://arxiv.org/abs/2101.08596}
    \item Schluter, J. and Gutenbrunner, G. \textit{EfficientLEAF: A Faster
    LEarnable Audio Frontend of Questionable Use}. arXiv:2207.05508.
    \url{https://arxiv.org/abs/2207.05508}
\end{itemize}

\end{document}
